\section{Methodology}
\label{sec:method}

In this section, we detail the mathematical and algorithmic formulation of our proposed \textbf{Task-Correlation Prior Regularization (TCPR)} framework. We first formalize the dynamic model merging paradigm, followed by the architecture of our Softmax-free Bounded Sigmoidal Router (BSigmoid-Router), and finally, the formulation of our TCPR regularizer.

\subsection{Dynamic Model Merging Formulation}
Let $W_{\text{base}}$ denote the weights of a pretrained base neural network, and let $W_k$ for $k \in \{1, \dots, K\}$ represent the weights of $K$ distinct task-specific experts fine-tuned from the same base model $W_{\text{base}}$. Following the formulation of Task Arithmetic \cite{ilharco2022editing}, we define the task vector $V_k$ for expert $k$ as:
\begin{equation}
    V_k = W_k - W_{\text{base}}
\end{equation}
In a multi-task setting, our objective is to combine these task vectors to form a single, unified set of weights $W_{\text{merged}}$. In dynamic model merging, instead of using static global coefficients, we compute input-dependent merging coefficients. For a given input batch $x$, the merged weights at each layer or block $l \in \{1, \dots, L\}$ are dynamically interpolated as:
\begin{equation}
    W_{\text{merged}}^{(l)} = W_{\text{base}}^{(l)} + \sum_{k=1}^K \bar{\alpha}_k V_k^{(l)}
\end{equation}
where $\bar{\alpha}_k \in [0, \lambda_{\text{max}}]$ represents the dynamic merging coefficient for task $k$, and $\lambda_{\text{max}}$ is a predefined scaling ceiling that prevents weight divergence.

\subsection{Bounded Sigmoidal Router (BSigmoid-Router)}
The merging coefficients $\bar{\alpha}_k$ are generated by a lightweight, trainable routing head. For an input batch $x \in \mathbb{R}^{B \times C \times H \times W}$ passed to a Vision Transformer backbone, we extract the patch embeddings and spatially average them to form a global representation vector $z(x) \in \mathbb{R}^{B \times D}$, where $D$ is the hidden dimension of the network.

The global representation is passed through a linear projection layer to compute raw routing logits $o(x) \in \mathbb{R}^{B \times K}$:
\begin{equation}
    o(x) = z(x) W_{\text{route}} + b_{\text{route}}
\end{equation}
where $W_{\text{route}} \in \mathbb{R}^{D \times K}$ is the routing projection weight matrix, and $b_{\text{route}} \in \mathbb{R}^K$ is the bias vector. Let $\mathbf{w}_k \in \mathbb{R}^D$ denote the $k$-th column of $W_{\text{route}}$, corresponding to the routing signature vector for task $k$.

Standard dynamic routers typically employ a Softmax activation to normalize the coefficients, which imposes a strict zero-sum competitive constraint: increasing the coefficient for one task necessarily decreases the coefficients for others. This competitive bottleneck is highly problematic in multi-task scenarios where multiple domains may be simultaneously relevant. 

To overcome this, we employ a decoupled, Softmax-free activation using independent sigmoid functions. We define the sample-level routing coefficients $\alpha_{k, b}$ for each sample $b$ in the batch as:
\begin{equation}
    \alpha_{k, b} = \lambda_{\text{max}} \times \text{Sigmoid}(o(x)_{b, k})
\end{equation}
where we fix the scale ceiling $\lambda_{\text{max}} = 0.3$. These sample-level coefficients are then averaged across the batch to yield the batch merging coefficients:
\begin{equation}
    \bar{\alpha}_k = \frac{1}{B} \sum_{b=1}^B \alpha_{k, b}
\end{equation}
By decoupling the activation pathways, our BSigmoid-Router allows multiple experts to be activated simultaneously or deactivated entirely based on the input's characteristics.

\subsection{Task-Correlation Prior Regularization (TCPR)}
During the low-data calibration phase (e.g., $N=16$ samples per task), the routing head parameters $W_{\text{route}}$ and $b_{\text{route}}$ are optimized using backpropagation. In this regime, the unconstrained optimization landscape is highly susceptible to overfitting, often resulting in representational collapse where the routing head completely shuts down the pathways for difficult or conflicting tasks.

To guide the optimization and prevent collapse, we introduce **Task-Correlation Prior Regularization (TCPR)**, which injects a pre-computed cross-task similarity matrix $S \in \mathbb{R}^{K \times K}$ as a prior. We define two distinct variants of the similarity matrix $S$:

\subsubsection{Parameter-Space Similarity Prior (TCPR-Param)}
This variant measures the structural similarity of the specialized experts in parameter space. We compute the cosine similarity of the task vectors $V_k^{(l)}$ across all layers $l \in \{1, \dots, L\}$ of the network:
\begin{equation}
    S^{\text{param}}_{i, j} = \frac{1}{L} \sum_{l=1}^L \frac{\langle V_i^{(l)}, V_j^{(l)} \rangle}{\|V_i^{(l)}\|_F \|V_j^{(l)}\|_F}
\end{equation}
where $\langle \cdot, \cdot \rangle$ represents the flat inner product and $\|\cdot\|_F$ denotes the Frobenius norm.

\subsubsection{Representation-Space Similarity Prior (TCPR-Rep)}
This variant measures the behavioral similarity of the experts in representation space on a generic validation subset. Let $Z_i, Z_j \in \mathbb{R}^{M \times D}$ be the intermediate representations (activations) extracted from the base model evaluated on $M$ validation samples:
\begin{equation}
    S^{\text{rep}}_{i, j} = \frac{\sum_{m=1}^M \langle Z_{i, m}, Z_{j, m} \rangle}{\sqrt{\sum_{m=1}^M \|Z_{i, m}\|^2_2} \sqrt{\sum_{m=1}^M \|Z_{j, m}\|^2_2}}
\end{equation}

\subsubsection{The TCPR Penalty Loss}
To resolve the scaling and collinearity issues inherent in unnormalized, positive-only prior matrices, we propose a mathematically rigorous reformulation of the TCPR penalty:

1. \textbf{Prior Matrix Centering}: Since similarity matrices $S$ computed via weights or activations are typically positive-only, we center the off-diagonal similarity coefficients by subtracting their mean:
\begin{equation}
    S^{\text{centered}}_{i, j} = \begin{cases}
        S_{i, j} - \mu_{\text{off}} & \text{if } i \neq j \\
        1.0 & \text{if } i = j
    \end{cases}
\end{equation}
where $\mu_{\text{off}} = \frac{1}{K(K-1)} \sum_{p \neq q} S_{p, q}$ is the mean of the off-diagonal similarity elements. This centering guarantees that tasks with above-average compatibility have positive priors ($S^{\text{centered}}_{i, j} > 0$), while tasks with below-average compatibility (conflicts) have negative priors ($S^{\text{centered}}_{i, j} < 0$).

2. \textbf{Signature Normalization}: To prevent weight explosion and resolve the dying gradient problem under small initialization scales, we compute the cosine similarity of the routing signatures by projecting them to unit spheres:
\begin{equation}
    \cos(\mathbf{w}_i, \mathbf{w}_j) = \frac{\mathbf{w}_i^T \mathbf{w}_j}{\|\mathbf{w}_i\|_2 \|\mathbf{w}_j\|_2}
\end{equation}
This forces the prior terms to remain strictly bounded in $[-1, 1]$.

The TCPR penalty is then defined as:
\begin{equation}
    \mathcal{L}_{\text{prior}}(W_{\text{route}}) = - \sum_{i=1}^K \sum_{j \ne i}^K S^{\text{centered}}_{i, j} \cos(\mathbf{w}_i, \mathbf{w}_j)
\end{equation}

The mechanisms of TCPR can be understood through its effect on task relationships:
\begin{itemize}
    \item \textbf{Positive Task Correlation ($S^{\text{centered}}_{i, j} > 0$):} For highly similar tasks, TCPR encourages their routing signatures to be positively aligned ($\cos(\mathbf{w}_i, \mathbf{w}_j) \to +1$), allowing the router to use similar activation cues to activate both experts cooperatively.
    \item \textbf{Negative Task Correlation ($S^{\text{centered}}_{i, j} < 0$):} For conflicting tasks, TCPR forces their routing signatures to be negatively aligned ($\cos(\mathbf{w}_i, \mathbf{w}_j) \to -1$), ensuring that the routing pathways remain decoupled, which completely prevents catastrophic interference.
\end{itemize}

The complete calibration objective minimized over the tiny calibration dataset is:
\begin{equation}
    \mathcal{L}_{\text{total}} = \mathcal{L}_{\text{CE}} + \beta \mathcal{L}_{\text{prior}}(W_{\text{route}}) + \gamma \|W_{\text{route}}\|_F^2
\end{equation}
which is equivalent to:
\begin{equation}
    \mathcal{L}_{\text{total}} = \mathcal{L}_{\text{CE}} - \beta \sum_{i=1}^K \sum_{j \ne i}^K S^{\text{centered}}_{i, j} \cos(\mathbf{w}_i, \mathbf{w}_j) + \gamma \|W_{\text{route}}\|_F^2
\end{equation}
where $\mathcal{L}_{\text{CE}}$ is the joint multi-task cross-entropy loss, $\beta$ is the regularization scaling strength, and $\gamma$ is the isotropic L2 weight decay factor (fixed at $10^{-4}$).
